\documentclass[10pt,aspectratio=169]{beamer}

% ------------------------------------------------------------
% Aula prática: Git, GitHub e RStudio
% Compilação: pdfLaTeX
% ------------------------------------------------------------
\usepackage[utf8]{inputenc}
\usepackage[T1]{fontenc}
\usepackage[brazil]{babel}
\usepackage{lmodern}
\usepackage{hyperref}
\usepackage{xcolor}
\usepackage{listings}
\usepackage{booktabs}

\usetheme{Madrid}
\usecolortheme{default}
\setbeamertemplate{navigation symbols}{}
\setbeamertemplate{footline}[frame number]

\definecolor{UFVBlue}{HTML}{003A70}
\definecolor{UFVGreen}{HTML}{2E7D32}
\definecolor{CodeBg}{HTML}{F6F8FA}
\definecolor{CodeFrame}{HTML}{D0D7DE}

\setbeamercolor{title}{fg=white,bg=UFVBlue}
\setbeamercolor{frametitle}{fg=white,bg=UFVBlue}
\setbeamercolor{structure}{fg=UFVBlue}
\setbeamercolor{block title}{fg=white,bg=UFVBlue}
\setbeamercolor{block body}{bg=gray!5,fg=black}

\hypersetup{
  colorlinks=true,
  linkcolor=UFVBlue,
  urlcolor=UFVBlue
}

\lstdefinestyle{terminal}{
  backgroundcolor=\color{CodeBg},
  frame=single,
  rulecolor=\color{CodeFrame},
  basicstyle=\ttfamily\scriptsize,
  breaklines=true,
  columns=fullflexible,
  keepspaces=true,
  showstringspaces=false
}
\lstset{style=terminal}

\title[Git, GitHub e RStudio]{Conexão e Controle de Versão\\com Git, GitHub e RStudio}
\subtitle{Aula prática, objetiva e mão na massa}
\author{EST 100 / EST 150 -- Ciência de Dados}
\institute{Universidade Federal de Viçosa}
\date{\today}

\begin{document}

% ------------------------------------------------------------
\begin{frame}
  \titlepage
\end{frame}

% ------------------------------------------------------------
\begin{frame}{Objetivo da aula}
\large
Ao final, cada estudante deverá conseguir:

\vspace{0.4cm}
\begin{itemize}
  \item instalar e configurar o \textbf{Git};
  \item conectar \textbf{GitHub} ao \textbf{RStudio};
  \item criar/clonar um projeto versionado;
  \item fazer o ciclo: \textbf{Pull} $\rightarrow$ editar $\rightarrow$ \textbf{Commit} $\rightarrow$ \textbf{Push}.
\end{itemize}

\vspace{0.4cm}
\begin{block}{Ideia central}
Controle de versão = histórico + segurança + colaboração + reprodutibilidade.
\end{block}
\end{frame}

% ------------------------------------------------------------
\begin{frame}{Instalar o necessário}
\large
\begin{itemize}
  \item \textbf{Git}: \href{https://git-scm.com/downloads}{git-scm.com/downloads}
  \item \textbf{GitHub}: \href{https://github.com/}{github.com}
  \item \textbf{RStudio Desktop}: \href{https://posit.co/download/rstudio-desktop/}{posit.co/download/rstudio-desktop}
\end{itemize}

\vspace{0.35cm}
\begin{block}{No Windows}
Durante a instalação do Git, pode manter as opções padrão.\newline
Depois, use o terminal \textbf{Git Bash} ou o terminal do RStudio.
\end{block}
\end{frame}

% ------------------------------------------------------------
\begin{frame}[fragile]{Configuração única do Git}
No terminal do RStudio: \textbf{Tools} $\rightarrow$ \textbf{Terminal} $\rightarrow$ \textbf{New Terminal}

\vspace{0.25cm}
\begin{lstlisting}
git --version

git config --global user.name "Seu Nome"
git config --global user.email "seu_email@exemplo.com"

git config --global init.defaultBranch main
git config --global core.autocrlf true

git config --global --list
\end{lstlisting}

\vspace{0.15cm}
\begin{itemize}
  \item Use o mesmo e-mail cadastrado no GitHub.
  \item \texttt{main} será o nome padrão do ramo principal.
\end{itemize}
\end{frame}

% ------------------------------------------------------------
\begin{frame}[fragile]{Conectar GitHub por SSH}
\textbf{1. Gerar chave SSH no terminal:}

\begin{lstlisting}
ssh-keygen -t ed25519 -C "seu_email@exemplo.com"
\end{lstlisting}

\textbf{2. Copiar a chave pública:}

\begin{lstlisting}
cat ~/.ssh/id_ed25519.pub
\end{lstlisting}

\textbf{3. Colar no GitHub:}

\begin{itemize}
  \item GitHub $\rightarrow$ \textbf{Settings}
  \item \textbf{SSH and GPG keys}
  \item \textbf{New SSH key}
\end{itemize}

\textbf{4. Testar:}

\begin{lstlisting}
ssh -T git@github.com
\end{lstlisting}
\end{frame}

% ------------------------------------------------------------
\begin{frame}{Criar repositório no GitHub}
\large
No GitHub:

\vspace{0.25cm}
\begin{enumerate}
  \item \textbf{New repository}
  \item Nome: \texttt{est100-primeiro-projeto}
  \item Marque: \textbf{Add a README file}
  \item Adicione: \texttt{.gitignore} $\rightarrow$ \textbf{R}
  \item License: opcional
  \item \textbf{Create repository}
\end{enumerate}

\vspace{0.35cm}
\begin{block}{Regra simples}
Um projeto de análise de dados $=$ uma pasta $=$ um repositório.
\end{block}
\end{frame}

% ------------------------------------------------------------
\begin{frame}{Clonar no RStudio}
\large
No RStudio:

\vspace{0.25cm}
\begin{enumerate}
  \item \textbf{File} $\rightarrow$ \textbf{New Project}
  \item \textbf{Version Control}
  \item \textbf{Git}
  \item Cole a URL SSH do GitHub:
\end{enumerate}

\vspace{0.15cm}
\begin{center}
\texttt{git@github.com:usuario/repositorio.git}
\end{center}

\vspace{0.25cm}
\begin{block}{Resultado esperado}
O RStudio abre um projeto com a aba \textbf{Git} visível.
\end{block}
\end{frame}

% ------------------------------------------------------------
\begin{frame}{Fluxo diário de trabalho}
\Large
\begin{center}
\textbf{Pull} $\rightarrow$ editar $\rightarrow$ salvar $\rightarrow$ \textbf{Stage} $\rightarrow$ \textbf{Commit} $\rightarrow$ \textbf{Push}
\end{center}

\vspace{0.4cm}
\large
Na aba \textbf{Git} do RStudio:

\vspace{0.2cm}
\begin{itemize}
  \item \textbf{Pull}: baixa mudanças do GitHub.
  \item \textbf{Stage}: escolhe o que entrará no registro.
  \item \textbf{Commit}: grava uma versão local.
  \item \textbf{Push}: envia para o GitHub.
\end{itemize}
\end{frame}

% ------------------------------------------------------------
\begin{frame}[fragile]{Comandos essenciais}
\begin{lstlisting}
# Ver situacao atual
git status

# Baixar atualizacoes
git pull

# Adicionar arquivos
git add .

# Registrar versao
git commit -m "mensagem curta e clara"

# Enviar ao GitHub
git push

# Ver historico
git log --oneline
\end{lstlisting}

\vspace{0.1cm}
\begin{block}{Boa mensagem de commit}
\texttt{corrige leitura dos dados} \quad ou \quad \texttt{adiciona grafico de dispersao}
\end{block}
\end{frame}

% ------------------------------------------------------------
\begin{frame}{Boas práticas mínimas}
\large
\begin{itemize}
  \item Faça \textbf{commit pequeno}: uma mudança por vez.
  \item Nunca versione: senhas, tokens, CPF, dados sensíveis.
  \item Use \texttt{.gitignore} para arquivos temporários.
  \item Antes de começar: \textbf{Pull}.
  \item Ao terminar: \textbf{Commit + Push}.
  \item README deve explicar: objetivo, dados, scripts e saída.
\end{itemize}
\end{frame}

% ------------------------------------------------------------
\begin{frame}{Problemas comuns}
\large
\begin{tabular}{@{}p{0.36\textwidth}p{0.56\textwidth}@{}}
\toprule
\textbf{Problema} & \textbf{Solução rápida} \\
\midrule
Aba Git não aparece & Verificar se o projeto foi criado com Git \\
Git não encontrado & Reinstalar Git e reiniciar o RStudio \\
Erro de autenticação & Conferir SSH key no GitHub \\
Conflito de arquivo & Abrir, escolher versão correta, salvar e commitar \\
Commit sem e-mail & Conferir \texttt{git config --global user.email} \\
\bottomrule
\end{tabular}
\end{frame}

% ------------------------------------------------------------
\begin{frame}{Atividade prática em 20 minutos}
\large
\begin{enumerate}
  \item Criar repositório no GitHub.
  \item Clonar no RStudio.
  \item Criar \texttt{analise.R}.
  \item Escrever três linhas de código.
  \item Commit: \texttt{primeira analise em R}.
  \item Push.
  \item Conferir no GitHub.
\end{enumerate}

\vspace{0.35cm}
\begin{block}{Entrega}
Enviar apenas o link do repositório.
\end{block}
\end{frame}

% ------------------------------------------------------------
\begin{frame}{Links essenciais}
\large
\begin{itemize}
  \item Git: \href{https://git-scm.com/downloads}{git-scm.com/downloads}
  \item Livro Pro Git: \href{https://git-scm.com/book/en/v2}{git-scm.com/book/en/v2}
  \item GitHub Docs -- SSH: \href{https://docs.github.com/en/authentication/connecting-to-github-with-ssh}{docs.github.com/.../ssh}
  \item GitHub Docs -- Git básico: \href{https://docs.github.com/en/get-started/git-basics}{docs.github.com/.../git-basics}
  \item Posit/RStudio -- Version Control: \href{https://docs.posit.co/ide/user/ide/guide/tools/version-control.html}{docs.posit.co/.../version-control}
  \item Happy Git with R: \href{https://happygitwithr.com/}{happygitwithr.com}
\end{itemize}
\end{frame}

% ------------------------------------------------------------
\begin{frame}
\centering
\Huge Obrigado!

\vspace{0.5cm}
\Large Agora: abrir o RStudio e versionar o primeiro projeto.
\end{frame}

\end{document}
